\documentclass[aspectratio=169,11pt]{beamer}
\usetheme{Madrid}
\usecolortheme{default}
\usepackage[T1]{fontenc}
\usepackage{lmodern}
\usepackage{booktabs,tabularx,array}
\definecolor{iutblue}{RGB}{20,57,92}
\definecolor{iutteal}{RGB}{0,112,122}
\setbeamercolor{structure}{fg=iutblue}
\setbeamercolor{block title}{bg=iutblue,fg=white}
\setbeamercolor{block title example}{bg=iutteal,fg=white}
\setbeamertemplate{navigation symbols}{}
\setbeamertemplate{caption}[numbered]
\setbeamersize{text margin left=8mm,text margin right=8mm}
\newcolumntype{Y}{>{\raggedright\arraybackslash}X}
\newcommand{\source}[1]{\par\vfill{\tiny\color{gray}Source: Criterion 8 SAR, #1.}}
\newcommand{\coursehead}[2]{\small\textbf{Delivery:} #1\par\textbf{Assessment:} #2\par\medskip}
\NewDocumentEnvironment{cotable}{+b}{\footnotesize\renewcommand{\arraystretch}{1.22}\begin{tabularx}{\textwidth}{@{}p{7mm}Y p{27mm}p{20mm}@{}}\toprule\textbf{CO}&\textbf{Expected student work (summarized)}&\textbf{Mapped POs}&\textbf{Domains}\\\midrule #1\bottomrule\end{tabularx}}{}
\title[Criterion 8]{Criterion 8: Program Outcomes (POs) and Assessment}
% \subtitle{Sections 8.1 and 8.2: Outcome Statements, Alignment, and Course Design}
\author[Department of CSE]{Department of Computer Science and Engineering}
\institute[IUT]{Islamic University of Technology}
% \date{Source document: 24 June 2026}

\begin{document}

\begin{frame}
  \titlepage
  % \vspace{-5mm}
  % \begin{center}\scriptsize Based on \texttt{[2026-06-24] v2 SAR-CSE-criterion-8.pdf}, printed pp. 125--131.\end{center}
\end{frame}

\begin{frame}{8.1: Program outcome (PO) statements and their mapping with the PEOs}

Graduates of the B.Sc. in Computer Science and Engineering (CSE) program are expected to achieve the following listed Program Outcomes (POs) by the time of graduation.

% \begin{columns}[T]
% \column{0.48\textwidth}
% \begin{block}{8.1: }
% \begin{itemize}
% \item State the twelve Program Outcomes (POs) expected at graduation.
% \item Map each PO to a Program Educational Objective (PEO).
% \item Explain why each relationship supports the program's objectives.
% \end{itemize}
% \end{block}
% \column{0.48\textwidth}
% \begin{block}{8.2: Translate outcomes into course work}
% \begin{itemize}
% \item Develop and review measurable Course Outcomes (COs).
% \item Specify delivery, learning domains, and assessment tools.
% \item Link course-level learning to the relevant POs.
% \end{itemize}
% \end{block}
% \end{columns}
% \begin{exampleblock}{The alignment to explain}
% Program objectives $\longleftarrow$ graduate outcomes $\longleftarrow$ course outcomes $\longleftarrow$ learning activities and assessment evidence.
% \end{exampleblock}
% \source{Sections 8.1--8.2, pp. 125--131}
\end{frame}

% \section{8.1 Program Outcomes and PEO Alignment}
\begin{frame}{Program Outcomes}
\small
\renewcommand{\arraystretch}{1.35}
\begin{tabularx}{\textwidth}{@{}p{10mm}Y@{}}
\toprule\textbf{POs} & \textbf{Descriptions}\\\midrule
PO1 & Apply knowledge of mathematics, natural science, engineering fundamentals and an engineering specialization as specified in WK1 to WK4 respectively to the solution of complex engineering problems.\\
PO2 & Identify, formulate, research literature and analyse complex engineering problems reaching substantiated conclusions using first principles of mathematics, natural sciences and engineering sciences. (WK1 to WK4)\\
PO3 & Design solutions for complex engineering problems and design systems, components or processes that meet specified needs with appropriate consideration for public health and safety, cultural, societal, and environmental considerations. (WK5)\\
PO4 & Conduct investigations of complex problems using research-based knowledge (WK8) and research methods including design of experiments, analysis and interpretation of data, and synthesis of information to provide valid conclusions.\\

\bottomrule
\end{tabularx}
% \source{Table 8.1, p. 125; statements condensed for presentation}
\end{frame}

\begin{frame}{Continued}
\small
\renewcommand{\arraystretch}{1.35}
\begin{tabularx}{\textwidth}{@{}p{10mm}Y@{}}
\toprule
\textbf{POs} & \textbf{Descriptions} \\
\midrule

PO5 & Create, select and apply appropriate techniques, resources, and modern engineering and IT tools, including prediction and modelling, to complex engineering problems, with an under standing of the limitations. (WK6)\\
PO6 & Apply reasoning informed by contextual knowledge to assess societal, health, safety, legal and cultural issues and the consequent responsibilities relevant to professional engineering practice and solutions to complex engineering problems. (WK7)\\

PO7 & Understand and evaluate the sustainability and impact of professional engineering work in the solution of complex engineering problems in societal and environmental contexts. (WK7) \\

PO8 & Apply ethical principles and commit to professional ethics and responsibilities and norms of engineering practice. (WK7) \\

PO9 & Function effectively as an individual, and as a member or leader in diverse teams and in multi-disciplinary settings. \\


\bottomrule
\end{tabularx}
\end{frame}

\begin{frame}{Continued}
\small
\renewcommand{\arraystretch}{1.35}
\begin{tabularx}{\textwidth}{@{}p{10mm}Y@{}}
\toprule
\textbf{POs} & \textbf{Descriptions} \\
\midrule

PO10 & Communicate effectively on complex engineering activities with the engineering community and with society at large, such as being able to comprehend and write effective reports and design documentation, make effective presentations, and give and receive clear instructions. \\

PO11 & Demonstrate knowledge and understanding of engineering management principles and economic decision-making and apply these to one's own work, as a member and leader in a team, to manage projects and in multidisciplinary environments. \\

PO12 & Recognize the need for, and have the preparation and ability to engage in independent and life-long learning in the broadest context of technological change. \\
\bottomrule
\end{tabularx}
\end{frame}

% \begin{frame}{8.1: Professional capabilities expected at graduation}
% \footnotesize
% \renewcommand{\arraystretch}{1.18}
% \begin{tabularx}{\textwidth}{@{}p{10mm}p{37mm}Y@{}}
% \toprule\textbf{POs} & \textbf{Descriptions}\\\midrule
% PO6&The engineer and society&Assess societal, health, safety, legal, and cultural issues and professional responsibilities.\\
% PO7&Environment and sustainability&Evaluate the societal and environmental impact and sustainability of engineering work.\\
% PO8&Ethics&Apply ethical principles and uphold professional responsibilities and norms.\\
% PO9&Individual work and teamwork&Work effectively independently and as a member or leader of diverse teams.\\
% PO10&Communication&Write reports and design documents; present effectively; give and receive clear instructions.\\
% PO11&Project management and finance&Apply engineering management and economic decision-making in individual and team projects.\\
% PO12&Life-long learning&Recognize the need for and undertake independent, continuing learning as technology changes.\\
% \bottomrule
% \end{tabularx}
% \source{Table 8.1, pp. 125--126; statements condensed for presentation}
% \end{frame}

% \begin{frame}{8.1: How the twelve POs support the four PEOs}
% \small
% \renewcommand{\arraystretch}{1.15}
% \begin{tabularx}{\textwidth}{@{}p{14mm}p{28mm}Y@{}}
% \toprule\textbf{PEO}&\textbf{Mapped POs}&\textbf{Reason for the alignment}\\\midrule
% PEO1&PO1--PO5&Computing knowledge, analysis, design, research, and modern tools support engineering solution development.\\
% PEO2&PO6--PO8&Societal responsibility, sustainability, and ethics support responsible professional practice.\\
% PEO3&PO12&Independent learning supports continuing professional development.\\
% PEO4&PO9--PO11&Teamwork, communication, and management support collaboration in diverse and multidisciplinary settings.\\
% \bottomrule
% \end{tabularx}
% \begin{block}{What the mapping documents}
% Every PO has a stated connection to a PEO; the report supplies a justification for each connection. The PO approval document is referenced in Section G.1.
% \end{block}
% \source{Table 8.2 and mapping justifications, pp. 126--127; PEO descriptions above summarize the rationale, not official PEO wording}
% \end{frame}

\section{8.2 Course Outcomes}
\begin{frame}{8.2: How course outcomes are developed and reviewed}
\small
\begin{enumerate}
\item \textbf{Instructors draft the course package:} outline, contents, assessment plan, and CO--PO mappings using the approved OBE curriculum and syllabus.
\item \textbf{Course coordinators harmonize sections:} align contents, assessment, and mappings across multiple sections of a course.
\item \textbf{The Departmental Committee reviews:} examine proposed materials, previous attainment, CQI findings, and instructor, student, alumni, and employer feedback.
\item \textbf{The Syllabus Committee examines curricular changes:} check alignment with the mission, POs, and PEOs; significant changes follow institutional approval through the Academic Council.
\item \textbf{The department finalizes the COs:} incorporate approved improvements into outcomes, outlines, and assessment strategies.
\end{enumerate}
\begin{exampleblock}{The recurring work}
Use evidence and feedback from earlier offerings to improve the next course design.
\end{exampleblock}
\source{Section 8.2, p. 127}
\end{frame}

\begin{frame}{8.2.1: What is documented for each course outcome?}
\begin{columns}[T]
\column{0.49\textwidth}
\begin{block}{The CO specification}
\begin{itemize}
\item A specific, measurable statement of expected knowledge, skills, or competence.
\item One or more mapped POs.
\item Learning-domain codes.
\item Delivery methods and learning activities.
\item Assessment tools that can demonstrate the outcome.
\end{itemize}
\end{block}
\column{0.47\textwidth}
\begin{exampleblock}{Example: AI theory, CO3}
\textbf{Student work:} recommend solutions to real-life AI problems using existing techniques.\par\medskip
\textbf{Alignment:} PO3 (design).\par
\textbf{Domain:} C6.\par
\textbf{Delivery:} lecture and demonstration.\par
\textbf{Evidence:} mid-term and final examinations.
\end{exampleblock}
\end{columns}
\medskip\scriptsize Domain families: C = cognitive; A = affective; P = psychomotor. Numeric codes on the following slides are reproduced from Table 8.3.
\source{Section 8.2.1 and Table 8.3, pp. 127--130}
\end{frame}

\begin{frame}{HUM 4441: Engineering Ethics}
\coursehead{Lecture, group discussion, and problem solving.}{Mid-term and final examinations.}
\begin{cotable}
CO1&Understand moral and ethical principles governing engineering practice and decisions.&PO6&C2\\
CO2&Examine social and cultural implications, professional responsibility, and accountability.&PO8&C4\\
CO3&Develop a framework for ethical analysis and engineering decision-making.&PO8&C3\\
CO4&Analyze and apply ethical implications of data technologies, privacy, security, and ownership.&PO6&C4\\
CO5&Understand how data-driven technologies create or perpetuate bias and discrimination.&PO6&C2\\
\end{cotable}
\source{Table 8.3, p. 128; CO descriptions summarized}
\end{frame}

\begin{frame}{CSE 4513: Software Engineering and OOD}
\coursehead{Lecture and problem solving.}{Mid-term and final examinations.}
\begin{cotable}
CO1&Explain software development fundamentals and requirement engineering.&PO1, PO2&C2\\
CO2&Apply design knowledge to model large software systems.&PO2&C3\\
CO3&Use techniques and tools to release and maintain quality software.&PO1, PO2&C3\\
\end{cotable}
\begin{exampleblock}{Work to explain}
Students move from understanding the development process to modeling systems and applying software release and maintenance practices.
\end{exampleblock}
\source{Table 8.3, p. 128; CO descriptions summarized}
\end{frame}

\begin{frame}{CSE 4711: Artificial Intelligence}
\coursehead{Lecture and demonstration.}{Mid-term and final examinations.}
\begin{cotable}
CO1&Apply general AI techniques to common problem types.&PO1&C3\\
CO2&Formulate problem scenarios as environments and agents, with justification.&PO2&C5\\
CO3&Recommend solutions to real-life AI problems using existing techniques.&PO3&C6\\
\end{cotable}
\begin{exampleblock}{Work to explain}
Apply AI knowledge $\longrightarrow$ formulate the problem $\longrightarrow$ propose a suitable solution.
\end{exampleblock}
\source{Table 8.3, p. 128; CO descriptions summarized}
\end{frame}

\begin{frame}{CSE 4712: Artificial Intelligence Lab}
\coursehead{Group discussion, demonstration, and problem solving.}{Lab viva, lab tasks, and lab reports.}
\begin{cotable}
CO1&Implement and analyze algorithms for complex navigation and strategy problems.&PO5&C3, C4\\
CO2&Formulate complex real-world scenarios as specific AI paradigms.&PO2&C5\\
CO3&Design, test, and evaluate intelligent agents under uncertainty using modern tools and frameworks.&PO4&C6, P5\\
CO4&Communicate analysis, design choices, and performance through structured technical reports.&PO10&P7, C5\\
\end{cotable}
\medskip\scriptsize The PO mappings here follow Table 8.3; Table 8.4 differs for CO1 and CO3 (see the review slide).
\source{Table 8.3, pp. 128--129; CO descriptions summarized}
\end{frame}

\begin{frame}{CSE 4610: Design Project}
\coursehead{Group discussion, demonstration, and project work.}{Presentation and project.}
\begin{cotable}
CO1&Identify research problems and implement innovative solutions with current technologies, tools, and algorithms.&PO1--PO5&C4\\
CO2&Demonstrate teamwork, ethics, responsibility, and project management to deliver a software system within constraints.&PO8, PO9, PO11&C3, A4\\
CO3&Design and evaluate computing solutions considering ethical, societal, and environmental implications.&PO6, PO7, PO8&C5, A3\\
CO4&Communicate outcomes through a structured research report and formal presentation.&PO10&C3, A3\\
CO5&Build a foundation for future research in the same or related domains.&PO12&C2, A2\\
\end{cotable}
\source{Table 8.3, p. 129; CO descriptions summarized}
\end{frame}

\begin{frame}{CSE 4700: Project/Thesis}
\coursehead{Discussion, demonstration, problem solving, and report writing.}{Presentation and report.}
\begin{cotable}
CO1&Identify and formulate a significant research problem through literature review and analysis of existing solutions.&PO1, PO2, PO4&C4, A4\\
CO2&Propose methodologies, models, or systems for the identified problem.&PO3&C6, A5, P5\\
CO3&Conduct experiments, simulations, or data-driven investigations; evaluate solutions and interpret results.&PO4, PO5&C4, A4, P4\\
CO5&Prepare a technical report and present findings for technical and non-technical audiences.&PO10&C5, P5, A7\\
CO6&Work independently and collaboratively, apply project management, and pursue continuous learning.&PO9, PO11, PO12&C6, P5, A7\\
\end{cotable}
\medskip\scriptsize CO numbering follows the source: Table 8.3 does not list a CO4 for this course.
\source{Table 8.3, pp. 129--130; CO descriptions summarized}
\end{frame}

\begin{frame}{CSE 4800: Project/Thesis}
\coursehead{Discussion, demonstration, problem solving, and report writing.}{Presentation and report.}
\begin{cotable}
CO2&Propose methodologies, models, or systems for the identified problem.&PO3&C6, A5, P5\\
CO3&Conduct experiments, simulations, or data-driven investigations; evaluate solutions and interpret results.&PO4, PO5&C4, A4, P4\\
CO4&Demonstrate ethical responsibility and professionalism in research, considering society and the environment.&PO6, PO7, PO8&C6, P6\\
CO5&Prepare a technical report and present findings for technical and non-technical audiences.&PO10&C5, P5, A7\\
CO6&Work independently and collaboratively, apply project management, and pursue continuous learning.&PO9, PO11, PO12&C6, P5, A7\\
\end{cotable}
\medskip\scriptsize CO numbering follows the source: Table 8.3 does not list a CO1 for this course.
\source{Table 8.3, p. 130; CO descriptions summarized}
\end{frame}

\begin{frame}{8.2.2: How the courses collectively contribute to POs}
\footnotesize
\renewcommand{\arraystretch}{1.1}
\begin{tabularx}{\textwidth}{@{}p{48mm}Y@{}}
\toprule\textbf{Course}&\textbf{PO coverage derived from the CO statements in Table 8.3}\\\midrule
HUM 4441: Engineering Ethics&PO6, PO8\\
CSE 4513: Software Engineering and OOD&PO1, PO2\\
CSE 4711: Artificial Intelligence&PO1, PO2, PO3\\
CSE 4712: Artificial Intelligence Lab&PO2, PO4, PO5, PO10\\
CSE 4610: Design Project&PO1--PO12 (all twelve)\\
CSE 4700: Project/Thesis&PO1, PO2, PO3, PO4, PO5, PO9, PO10, PO11, PO12\\
CSE 4800: Project/Thesis&PO3, PO4, PO5, PO6, PO7, PO8, PO9, PO10, PO11, PO12\\
\bottomrule
\end{tabularx}
\begin{exampleblock}{Why the distribution matters}
Theory develops knowledge and analysis; laboratory work adds implementation and investigation; projects integrate technical and professional capabilities.
\end{exampleblock}
\source{Section 8.2.2 and Tables 8.3--8.4, pp. 128--131; coverage uses Table 8.3 where the tables differ}
\end{frame}

\begin{frame}{Source review: reconcile the CO--PO tables}
\footnotesize
Table 8.3 and Table 8.4 contain differences that should be checked before finalizing the SAR or presentation.
\par\medskip\noindent
\renewcommand{\arraystretch}{1.4}
\begin{tabularx}{\textwidth}{@{}p{36mm}Y Y@{}}
\toprule\textbf{Course / CO}&\textbf{Table 8.3: statements}&\textbf{Table 8.4: mapping matrix}\\\midrule
CSE 4712 / CO1&PO5&PO4\\
CSE 4712 / CO3&PO4&PO3\\
CSE 4700 / CO5&PO10&PO6\\
CSE 4800 / CO5&PO10&No PO marked\\
\bottomrule
\end{tabularx}
\begin{block}{Treatment in this draft}
Course slides and the coverage summary consistently use Table 8.3. Retain the original CO numbering and learning-domain codes pending confirmation against the approved course outlines.
\end{block}
\source{Table 8.3, pp. 128--130; Table 8.4, p. 131}
\end{frame}

\begin{frame}{What Sections 8.1 and 8.2 document}
\begin{enumerate}
\item \textbf{Graduate expectations:} twelve defined POs, with an approval reference.
\item \textbf{Program alignment:} PO--PEO mappings and a rationale for every relationship.
\item \textbf{Course design:} CO statements, learning domains, delivery methods, and assessment tools for seven selected core courses.
\item \textbf{Review responsibilities:} instructor preparation, coordinator harmonization, departmental review, and formal approval of significant changes.
\item \textbf{Traceability:} each course outcome connects student work to one or more graduate outcomes.
\end{enumerate}
\begin{exampleblock}{Presentation takeaway}
These sections explain how the program defines its outcomes and organizes course work to support them.
\end{exampleblock}
\source{Sections 8.1--8.2, pp. 125--131}
\end{frame}

\end{document}
